\documentclass{amsart}
% For dealing with references we use the comment environment
\usepackage{verbatim}
\newenvironment{reference}{\comment}{\endcomment}
%\newenvironment{reference}{}{}
\newenvironment{slogan}{\comment}{\endcomment}
\newenvironment{history}{\comment}{\endcomment}

% For commutative diagrams we use Xy-pic
\usepackage[all]{xy}

% We use 2cell for 2-commutative diagrams.
\xyoption{2cell}
\UseAllTwocells

% We use multicol for the list of chapters between chapters
\usepackage{multicol}

% This is generally recommended for better output
\usepackage{lmodern}
\usepackage[T1]{fontenc}

% For cross-file-references

% Package for hypertext links:
\usepackage{hyperref}

% For any local file, say "hello.tex" you want to link to please

% Theorem environments.
%
\theoremstyle{plain}
\newtheorem{theorem}[subsection]{Theorem}
\newtheorem{proposition}[subsection]{Proposition}
\newtheorem{lemma}[subsection]{Lemma}

\theoremstyle{definition}
\newtheorem{definition}[subsection]{Definition}
\newtheorem{example}[subsection]{Example}
\newtheorem{exercise}[subsection]{Exercise}
\newtheorem{situation}[subsection]{Situation}

\theoremstyle{remark}
\newtheorem{remark}[subsection]{Remark}
\newtheorem{remarks}[subsection]{Remarks}

\numberwithin{equation}{subsection}

% Macros
%
\def\lim{\mathop{\mathrm{lim}}\nolimits}
\def\colim{\mathop{\mathrm{colim}}\nolimits}
\def\Spec{\mathop{\mathrm{Spec}}}
\def\Hom{\mathop{\mathrm{Hom}}\nolimits}
\def\Ext{\mathop{\mathrm{Ext}}\nolimits}
\def\SheafHom{\mathop{\mathcal{H}\!\mathit{om}}\nolimits}
\def\SheafExt{\mathop{\mathcal{E}\!\mathit{xt}}\nolimits}
\def\Sch{\mathit{Sch}}
\def\Mor{\mathop{\mathrm{Mor}}\nolimits}
\def\Ob{\mathop{\mathrm{Ob}}\nolimits}
\def\Sh{\mathop{\mathit{Sh}}\nolimits}
\def\NL{\mathop{N\!L}\nolimits}
\def\CH{\mathop{\mathrm{CH}}\nolimits}
\def\proetale{{pro\text{-}\acute{e}tale}}
\def\etale{{\acute{e}tale}}
\def\QCoh{\mathit{QCoh}}
\def\Ker{\mathop{\mathrm{Ker}}}
\def\Im{\mathop{\mathrm{Im}}}
\def\Coker{\mathop{\mathrm{Coker}}}
\def\Coim{\mathop{\mathrm{Coim}}}

% Boxtimes
%
\DeclareMathSymbol{\boxtimes}{\mathbin}{AMSa}{"02}

%
% Macros for moduli stacks/spaces
%
\def\QCohstack{\mathcal{QC}\!\mathit{oh}}
\def\Cohstack{\mathcal{C}\!\mathit{oh}}
\def\Spacesstack{\mathcal{S}\!\mathit{paces}}
\def\Quotfunctor{\mathrm{Quot}}
\def\Hilbfunctor{\mathrm{Hilb}}
\def\Curvesstack{\mathcal{C}\!\mathit{urves}}
\def\Polarizedstack{\mathcal{P}\!\mathit{olarized}}
\def\Complexesstack{\mathcal{C}\!\mathit{omplexes}}
% \Pic is the operator that assigns to X its picard group, usage \Pic(X)
% \Picardstack_{X/B} denotes the Picard stack of X over B
% \Picardfunctor_{X/B} denotes the Picard functor of X over B
\def\Pic{\mathop{\mathrm{Pic}}\nolimits}
\def\Picardstack{\mathcal{P}\!\mathit{ic}}
\def\Picardfunctor{\mathrm{Pic}}
\def\Deformationcategory{\mathcal{D}\!\mathit{ef}}

\setlength{\emergencystretch}{3em}
\begin{document}
\title[Stacks Project, Tag 0A9P]{253 --- Localization on the base of the right adjoint for proper pushforward}
\maketitle
\noindent Copyright (C) 2005--2025 Johan de Jong. Stacks Project authors. Source: \href{https://stacks.math.columbia.edu/tag/0A9P}{Tag 0A9P}.

\noindent Editorial difficulty note (not author text). Difficulty H2: The proof detects a hypothetical support violation with a perfect complex, lifts that detection across the open immersion, and annihilates its adjoint map by a perfect complex on the base. The projection formula then gives a contradiction while preserving the nonzero restricted map. The actual support control of the adjoint is a nonroutine argument..

\noindent Editorial provenance note (not author text). History: excerpt from revision \texttt{a0444\allowbreak 6e57e\allowbreak c1fbc\allowbreak 252a8\allowbreak 71afc\allowbreak ec775\allowbreak 2fb28\allowbreak 07b14}, duality.tex, lines 888--944. Reference presentation was adapted; original-author context and core support blocks were retained or added. The mathematical wording is unchanged. Licensed under GNU FDL 1.2 or later; no invariant sections or cover texts. See ../licenses/COPYING. Context blocks reproduce author definitions and supporting results; they are not additional counted problems.

\noindent
Let $f : X \to Y$ be a morphism of quasi-compact and quasi-separated
schemes. Let $V \subset Y$ be a quasi-compact open subscheme and set
$U = f^{-1}(V)$. This gives a cartesian square
$$
\xymatrix{
U \ar[r]_{j'} \ar[d]_{f|_U} & X \ar[d]^f \\
V \ar[r]^j & Y
}
$$
as in (\href{https://stacks.math.columbia.edu/tag/0A9J}{equation base change (Tag 0A9J)}). By Lemma \href{https://stacks.math.columbia.edu/tag/0A9K}{lemma flat precompose pus (Tag 0A9K)}
the map $\xi : a \circ Rj_* \leftarrow Rj'_* \circ a'$ is an isomorphism
where $a$ and $a'$ are the right adjoints to
$Rf_*$ and $R(f|_U)_*$. We obtain a transformation
of functors $D_\QCoh(\mathcal{O}_Y) \to D_\QCoh(\mathcal{O}_U)$
\begin{equation}
\label{equation-sheafy}
(j')^* \circ a \to
(j')^* \circ a \circ Rj_* \circ j^* \xrightarrow{\xi^{-1}}
(j')^* \circ Rj'_* \circ a' \circ j^* \to a' \circ j^*
\end{equation}
where the first arrow comes from $\text{id} \to Rj_* \circ j^*$
and the final arrow from the isomorphism $(j')^* \circ Rj'_* \to \text{id}$.
In particular, we see that (\ref{equation-sheafy}) is an isomorphism
when evaluated on $K$ if and only if $a(K)|_U \to a(Rj_*(K|_V))|_U$
is an isomorphism.

\begin{lemma}
\label{lemma-when-sheafy}
Let $f : X \to Y$ be a morphism of quasi-compact and quasi-separated
schemes. Let $a$ be the right adjoint to
$Rf_* : D_\QCoh(\mathcal{O}_X) \to D_\QCoh(\mathcal{O}_Y)$.
Let $V \subset Y$ be quasi-compact open with inverse image $U \subset X$.
\begin{enumerate}
\item For every $Q \in D_\QCoh^+(\mathcal{O}_Y)$
supported on $Y \setminus V$ the image $a(Q)$ is supported on
$X \setminus U$ if and only if (\ref{equation-sheafy})
is an isomorphism on all $K$ in $D_\QCoh^+(\mathcal{O}_Y)$.
\item For every $Q \in D_\QCoh(\mathcal{O}_Y)$
supported on $Y \setminus V$ the image $a(Q)$ is supported on
$X \setminus U$ if and only if (\ref{equation-sheafy})
is an isomorphism on all $K$ in $D_\QCoh(\mathcal{O}_Y)$.
\item If $a$ commutes with direct sums, then the equivalent conditions of
(1) imply the equivalent conditions of (2).
\end{enumerate}
\end{lemma}

\begin{lemma}
\label{lemma-proper-noetherian}
Let $Y$ be a quasi-compact and quasi-separated scheme.
Let $f : X \to Y$ be a proper morphism. If\footnote{This proof works for those
morphisms of quasi-compact and quasi-separated schemes such that
$Rf_*P$ is pseudo-coherent for all $P$ perfect on $X$. It follows
easily from a theorem of Kiehl \href{https://stacks.math.columbia.edu/bibliography}{Bibliography: Kiehl} that this holds if
$f$ is proper and pseudo-coherent. This is the correct generality
for this lemma and some of the other results in this chapter.}
\begin{enumerate}
\item $f$ is flat and of finite presentation, or
\item $Y$ is Noetherian
\end{enumerate}
then the equivalent conditions of Lemma \ref{lemma-when-sheafy} part (1)
hold for all quasi-compact opens $V$ of $Y$.
\end{lemma}

\begin{proof}
Let $Q \in D^+_\QCoh(\mathcal{O}_Y)$ be supported on $Y \setminus V$.
To get a contradiction, assume that $a(Q)$ is not supported on
$X \setminus U$. Then we can find a perfect complex $P_U$ on $U$
and a nonzero map $P_U \to a(Q)|_U$ (follows from
Derived Categories of Schemes, Theorem
\href{https://stacks.math.columbia.edu/tag/09IS}{theorem bondal van den Bergh (Tag 09IS)}). Then using
Derived Categories of Schemes, Lemma
\href{https://stacks.math.columbia.edu/tag/09IM}{lemma lift map from perfect complex with support (Tag 09IM)}
we may assume there is a perfect complex $P$ on $X$ and a map
$P \to a(Q)$ whose restriction to $U$ is nonzero.
By definition of $a$ this map
is adjoint to a map $Rf_*P \to Q$.

\medskip\noindent
The complex $Rf_*P$ is pseudo-coherent. In case (1) this follows
from Derived Categories of Schemes, Lemma
\href{https://stacks.math.columbia.edu/tag/0CSD}{lemma flat proper pseudo coherent direct image general (Tag 0CSD)}.
In case (2) this follows from
Derived Categories of Schemes, Lemmas
\href{https://stacks.math.columbia.edu/tag/08E2}{lemma direct image coherent (Tag 08E2)} and
\href{https://stacks.math.columbia.edu/tag/08E8}{lemma identify pseudo coherent noetherian (Tag 08E8)}.
Thus we may apply
Derived Categories of Schemes, Lemma
\href{https://stacks.math.columbia.edu/tag/0A9C}{lemma map from pseudo coherent to complex with support (Tag 0A9C)}
and get a map $I \to \mathcal{O}_Y$ of perfect complexes
whose restriction to $V$ is an isomorphism such that the composition
$I \otimes^\mathbf{L}_{\mathcal{O}_Y} Rf_*P \to Rf_*P \to Q$ is zero.
By Derived Categories of Schemes, Lemma
\href{https://stacks.math.columbia.edu/tag/08EU}{lemma cohomology base change (Tag 08EU)}
we have $I \otimes^\mathbf{L}_{\mathcal{O}_Y} Rf_*P =
Rf_*(Lf^*I \otimes^\mathbf{L}_{\mathcal{O}_X} P)$.
We conclude that the composition
$$
Lf^*I \otimes^\mathbf{L}_{\mathcal{O}_X} P \to P \to a(Q)
$$
is zero. However, the restriction to $U$ is the map
$P|_U \to a(Q)|_U$ which we assumed to be nonzero.
This contradiction finishes the proof.
\end{proof}
\end{document}
